% !TeX root = ../main.tex

\chapter{Conclusion and Future Work}

\section{Conclusion}

This thesis studies how PMDetector-style workflows can be evaluated as
audit-assistance systems rather than only as binary vulnerability classifiers.
Price manipulation vulnerabilities are difficult to detect because the
dangerous behavior depends on both data flow and economic context. A price-like
source is not always vulnerable, and an economically meaningful sink is not
always controlled by a manipulable value. Static analysis can identify possible
flows, but it tends to over-approximate. LLMs can reason about code context and
missing mitigations, but they require concrete evidence and structured output
to be useful.

The prototype developed for this thesis combines these two strengths. It uses
Slither-based taint analysis to identify candidate flows, canonical source
records to merge equivalent sources, representative path generation to reduce
raw path redundancy, and a staged LLM workflow to classify sources, triage sink
impact, and emit audit-style root-cause findings. The output is designed to be
traceable: each final finding refers back to concrete static-analysis evidence.

The main argument of this thesis is that protocol-level precision, recall, and
F1-score are insufficient for evaluating such systems. A detector can correctly
label a project as vulnerable while still pointing to the wrong path, naming
the wrong function, explaining the wrong mitigation, or producing many extra
reports. These failures matter because auditors do not only need a yes-or-no
answer. They need precise, actionable, and low-noise reports.

To address this gap, the thesis proposes a report-level evaluation framework
based on audit findings. The framework combines binary metrics with
localization checks, semantic explanation similarity, mitigation agreement, and
an over-reporting penalty. This design is intended to measure whether generated
reports are actually useful in an audit workflow.

\section{Current Status}

The current implementation should be understood as an initial prototype. It has
working components for taint analysis, source and sink tracking, path grouping,
LLM prompt generation, staged LLM execution, JSON validation, and preliminary
dataset organization. Existing artifacts show that the workflow can produce
plausible findings for projects such as Spartan and Predy. They also reveal
important limitations, such as missing sink coverage in some projects.

The current dataset work also remains in progress. The collected audit-finding
CSV contains useful metadata for filtering and evaluation, and the HIGH dataset
slice includes build-readiness checks. However, final quantitative evaluation
requires a clearly frozen build-verified subset and manually reviewed ground
truth labels.

\section{Future Work}

\subsection{Complete Sanity Verification}

The most important implementation task is to complete sanity verification.
Static reachability and LLM reasoning should be checked against concrete
mitigation patterns. Future rules should cover TWAP adequacy, oracle staleness,
deviation checks, slippage checks, access control, time-based controls,
fee-on-transfer behavior, and protocol-specific guard conditions. These checks
would reduce false positives and make LLM findings more defensible.

\subsection{Expand Sink Coverage}

Preliminary artifacts show that source identification can succeed while sink
matching fails. Future work should expand sink definitions beyond the current
patterns. In particular, the analyzer should better capture protocol-specific
accounting updates, vault share calculations, liquidity-position updates,
reward eligibility checks, and liquidation-adjacent helper functions. Each new
sink definition should be tested against real audit findings.

\subsection{Freeze a Build-Verified Dataset}

The final evaluation should use a frozen dataset split. Each project should
record the repository commit, compiler version, toolchain, build command,
required environment variables, and analysis target. Findings should be
categorized as price manipulation, related manipulation, or out of scope.
Without a frozen dataset, results will be hard to reproduce and compare.

\subsection{Calibrate Semantic Matching}

The semantic explanation score currently uses a working threshold. Future work
should build a small validation set of prediction-ground-truth pairs labeled by
human reviewers. The embedding model and threshold should then be calibrated
against this validation set. This step is necessary before semantic similarity
can be treated as a reliable metric.

\subsection{Implement Stronger Baselines}

A direct PMDetector reproduction is difficult without the original source code,
but useful baselines are still possible. One baseline can use the same static
analysis front end with a single PMDetector-style prompt that emits
\texttt{is\_high\_risk}, vulnerable functions, vulnerable groups, and a reason.
Another baseline can evaluate static analysis alone without LLM filtering.
Comparing these baselines with the staged workflow would clarify which parts of
the design improve report-level usefulness.

\subsection{Human Auditor Study}

The strongest validation would involve human auditors. Given a set of generated
reports, auditors could rate whether each report helps them locate the issue,
understand the root cause, and decide on mitigation. Such a study would
directly evaluate the thesis claim that report-level metrics better capture
audit usefulness than binary labels alone.

\section{Final Remarks}

LLM-assisted security tools should not be evaluated only by whether they can
say that a project is vulnerable. For practical auditing, the harder and more
important question is whether they can explain the vulnerability accurately,
localize it precisely, and avoid unnecessary noise. This thesis takes a step
toward that goal by implementing a PMDetector-inspired prototype and proposing
an audit-centered evaluation framework for price manipulation detection.
